\chapter{Holomorphic Functions}

\section{Holomorphic functions of several variables}

\subsection*{Notation and preliminaries}

Let the $n$-dimensional complex space be
\[
  \C^n = \{\, z = (z_1, \dots, z_n) \mid z_\nu \in \C \,\},
\]
and write each coordinate as $z_\nu = x_\nu + i y_\nu$
($x_\nu, y_\nu \in \R$, $i = \sqrt{-1}$). Through the correspondence
$z \mapsto (x_1, y_1, \dots, x_n, y_n)$ taking real and imaginary parts, we
identify $\C^n \cong \R^{2n}$.

Define the norm of $z = (z_1, \dots, z_n)$ by
\[
  |z| = \Bigl( \sum_{\nu=1}^{n} |z_\nu|^2 \Bigr)^{1/2}.
\]
For $z, w \in \C^n$ the following hold:
\begin{align}
  |z_\nu| &\le |z| \le \sum_{\nu=1}^n |z_\nu|, \\
  |z + w| &\le |z| + |w|, \\
  \bigl|\, |z| - |w| \,\bigr| &\le |z - w|.
\end{align}

\subsection*{Holomorphic functions}

\begin{definition}\label{def:holomorphic}
Let $D \subseteq \C^n$. We say that $f(z) = f(z_1, \dots, z_n)$ is a
\emph{holomorphic function} on $D$ if $f$ is continuous on $D$ and is
holomorphic in each variable $z_\nu$ separately (the other variables being held
fixed).
\end{definition}

\subsection*{Domains and polydiscs}

\begin{definition}[domain]
A connected open subset of $\C^n$ is called a \emph{domain}.
\end{definition}

\begin{note}
A topological space $X$ is \emph{connected} if every decomposition
$X = U \cup V$ with $U \cap V = \varnothing$ and $U, V$ open satisfies
$U = \varnothing$ or $V = \varnothing$ (that is, $X$ cannot be split into two
nonempty open sets). For open subsets of $\C^n$, connectedness is equivalent to
path-connectedness (any two points can be joined by a path within the set).
\end{note}

\begin{definition}[polydisc]
For $a = (a_1, \dots, a_n) \in \C^n$ and $r = (r_1, \dots, r_n)$ ($r_\nu > 0$),
\[
  P(a, r) = \{\, z \in \C^n \mid |z_\nu - a_\nu| < r_\nu,\ \nu = 1, \dots, n \,\}
\]
is called the \emph{polydisc} with center $a$ and radius $r$.
\end{definition}

\begin{note}
With the maximum norm $\|z\| = \max_{1 \le \nu \le n} |z_\nu|$, a polydisc is
the ``ball'' $\{\, z \mid \|z - a\| < r \,\}$ with respect to $\|z - a\|$ (when
$r_1 = \dots = r_n = r$).
\end{note}

\subsection*{Wirtinger derivatives (complex derivatives) and the Cauchy--Riemann equations}

For $z_\nu = x_\nu + i y_\nu$, define the \emph{Wirtinger derivatives} by
\[
  \frac{\partial}{\partial z_\nu}
    = \frac{1}{2}\Bigl( \frac{\partial}{\partial x_\nu} - i \frac{\partial}{\partial y_\nu} \Bigr),
  \qquad
  \frac{\partial}{\partial \bar z_\nu}
    = \frac{1}{2}\Bigl( \frac{\partial}{\partial x_\nu} + i \frac{\partial}{\partial y_\nu} \Bigr).
\]
In what follows we split a $C^1$ function $f$ into real and imaginary parts,
$f = u + i v$ ($u, v$ real-valued).

\begin{definition}[holomorphy, Cauchy--Riemann equations]
We say that $f = u + i v$ is \emph{holomorphic} in the variable $z_\nu$ if
\[
  \frac{\partial f}{\partial \bar z_\nu} = 0 .
\]
This is equivalent to the Cauchy--Riemann equations
\[
  \frac{\partial u}{\partial x_\nu} = \frac{\partial v}{\partial y_\nu}, \qquad
  \frac{\partial u}{\partial y_\nu} = -\frac{\partial v}{\partial x_\nu} .
\]
\end{definition}

\begin{proof}
By definition,
\[
  \frac{\partial f}{\partial \bar z_\nu}
  = \frac{1}{2}\Bigl( \frac{\partial}{\partial x_\nu} + i \frac{\partial}{\partial y_\nu} \Bigr)(u + i v)
  = \frac{1}{2}\Bigl[ \Bigl( \frac{\partial u}{\partial x_\nu} - \frac{\partial v}{\partial y_\nu} \Bigr)
      + i \Bigl( \frac{\partial v}{\partial x_\nu} + \frac{\partial u}{\partial y_\nu} \Bigr) \Bigr].
\]
A complex number is $0$ exactly when its real and imaginary parts both vanish,
so $\partial f / \partial \bar z_\nu = 0$ is equivalent to the Cauchy--Riemann
equations above.
\end{proof}

\subsection*{Convergence of power series}

\begin{note}
If, in a neighborhood of each point $a$ of $W$, $f(z)$ can be expanded as a
convergent power series
\[
  f(z) = \sum_{k_1, \dots, k_n \ge 0}
    c_{k_1 \cdots k_n} (z_1 - a_1)^{k_1} \cdots (z_n - a_n)^{k_n},
\]
then $f(z)$ is a holomorphic function.
\end{note}

\begin{proposition}\label{cl:abel}
If the power series $\sum_{k_1, \dots, k_n \ge 0} c_{k_1 \cdots k_n}
(z_1 - a_1)^{k_1} \cdots (z_n - a_n)^{k_n}$ converges at some point $w$, then it
converges absolutely for every $z$ satisfying
\[
  |z_\nu - a_\nu| < |w_\nu - a_\nu| \qquad (\nu = 1, \dots, n).
\]
\end{proposition}

\begin{proof}
We may assume $a = 0$ (replace $z_\nu - a_\nu$ by $z_\nu$). Since the series
converges at $w$, its terms are bounded: there is a constant $C < \infty$ with
\[
  |c_{k_1 \cdots k_n} \, w_1^{k_1} \cdots w_n^{k_n}| \le C
  \qquad (\forall\, k_1, \dots, k_n).
\]
For $|z_\nu| < |w_\nu|$,
\[
  |c_{k_1 \cdots k_n} \, z_1^{k_1} \cdots z_n^{k_n}|
  = |c_{k_1 \cdots k_n} \, w_1^{k_1} \cdots w_n^{k_n}|
    \prod_{\nu=1}^n \left| \frac{z_\nu}{w_\nu} \right|^{k_\nu}
  \le C \prod_{\nu=1}^n \left| \frac{z_\nu}{w_\nu} \right|^{k_\nu},
\]
and therefore
\[
  \sum_{k_1, \dots, k_n \ge 0} |c_{k_1 \cdots k_n} \, z_1^{k_1} \cdots z_n^{k_n}|
  \le C \prod_{\nu=1}^n \sum_{k=0}^\infty \left| \frac{z_\nu}{w_\nu} \right|^{k}
  = C \prod_{\nu=1}^n \frac{1}{1 - |z_\nu / w_\nu|} < \infty .
\]
Hence the series converges absolutely.
\end{proof}

\subsection*{Theorem 1: power-series expansion}

\begin{theorem}\label{thm:power-series}
Let $W \subseteq \C^n$ be a domain. If $f(z)$ is continuous on $W$ and
holomorphic in each variable $z_\nu$ separately (the other variables being held
fixed), then $f(z)$ is a holomorphic function of the $n$ variables
$z = (z_1, \dots, z_n)$.
\end{theorem}

\begin{proof}
Let $a$ be an arbitrary point of $W$. Since $W$ is open, we may choose
$r = (r_1, \dots, r_n)$ with
\[
  \{\, z \mid |z_\nu - a_\nu| \le r_\nu,\ \nu = 1, \dots, n \,\} \subset W .
\]
As $f$ is holomorphic in $z_1$, the Cauchy integral formula in the $z_1$-plane
gives
\[
  f(z_1, z_2, \dots, z_n)
  = \frac{1}{2\pi i} \oint_{|w_1 - a_1| = r_1}
      \frac{f(w_1, z_2, \dots, z_n)}{w_1 - z_1}\, dw_1 .
\]
Repeating this for each variable, and using that $f$ is continuous so that the
iterated integral equals the multiple integral (the order of integration may be
exchanged), we obtain
\[
  f(z) = \frac{1}{(2\pi i)^n}
    \oint_{|w_1 - a_1| = r_1} \!\!\!\! \cdots \oint_{|w_n - a_n| = r_n}
    \frac{f(w_1, \dots, w_n)}{(w_1 - z_1) \cdots (w_n - z_n)}\, dw_1 \cdots dw_n .
\]
Now, noting that $|z_\nu - a_\nu| < r_\nu = |w_\nu - a_\nu|$, each factor expands
as
\[
  \frac{1}{w_\nu - z_\nu}
  = \frac{1}{w_\nu - a_\nu} \cdot
    \frac{1}{1 - \dfrac{z_\nu - a_\nu}{w_\nu - a_\nu}}
  = \sum_{k=0}^\infty \frac{(z_\nu - a_\nu)^k}{(w_\nu - a_\nu)^{k+1}} .
\]
Substituting these and integrating term by term,
\[
  f(z) = \sum_{k_1, \dots, k_n \ge 0}
    c_{k_1 \cdots k_n} (z_1 - a_1)^{k_1} \cdots (z_n - a_n)^{k_n},
\]
where
\[
  c_{k_1 \cdots k_n} = \frac{1}{(2\pi i)^n}
    \oint_{|w_1 - a_1| = r_1} \!\!\!\! \cdots \oint_{|w_n - a_n| = r_n}
    \frac{f(w_1, \dots, w_n)}
         {(w_1 - a_1)^{k_1 + 1} \cdots (w_n - a_n)^{k_n + 1}}\, dw_1 \cdots dw_n .
\]
Hence $f$ is expanded as a convergent power series in a neighborhood of each
point of $W$, and is therefore a holomorphic function of $n$ variables.
\end{proof}
